% Generated draft: requires mathematical and source review.
\documentclass[11pt]{article}
\usepackage[margin=25mm]{geometry}
\usepackage{fontspec}
\setmainfont{Latin Modern Roman}
\usepackage{amsmath,amssymb,mathtools,unicode-math}
\setmathfont{Latin Modern Math}
\usepackage{enumitem,needspace,xurl,hyperref}
\setlength{\parindent}{0pt}
\setlength{\parskip}{5pt}
\setlength{\emergencystretch}{3em}
\newcommand{\N}{\mathbb{N}}
\newcommand{\Z}{\mathbb{Z}}
\newcommand{\Q}{\mathbb{Q}}
\newcommand{\R}{\mathbb{R}}
\newcommand{\C}{\mathbb{C}}
\newcommand{\F}{\mathbb{F}}
\newcommand{\A}{\mathbb{A}}
\newcommand{\PP}{\mathbb{P}}
\newcommand{\E}{\mathbb{E}}
\newcommand{\eps}{\varepsilon}
\newcommand{\dd}{\,\mathrm{d}}
\begin{document}
\section*{UnsolvedMath ID: 10400150}
\textbf{Source problem code:} catalogue entry
\subsection{Definitions and mathematical statement}
A spin topological quantum field theory is a topological quantum field theory whose bordisms and boundary objects carry spin structures, with the resulting assignments respecting composition by gluing and the monoidal operation of disjoint union. In the usual nonextended, vector-space-valued convention in dimension \(d\), this is expressed by a symmetric monoidal functor
\[
 Z:\operatorname{Bord}^{\mathrm{Spin}}_d\longrightarrow\operatorname{Vect}_k,
\]
where objects are closed spin \((d-1)\)-manifolds and morphisms are spin \(d\)-dimensional bordisms, up to the appropriate spin-preserving equivalence; \(k\) is a coefficient field. This describes a conventional framework only: the supplied excerpt fixes neither \(d\), \(k\), nor this target or extension convention.

\emph{Problem.} Construct other spin TQFTs, meaning examples additional to the examples intended in the preceding source discussion. The excerpt does not identify those existing examples or an equivalence criterion for deciding that a construction is new. Thus the class to be extended must be recovered before the word ``other'' determines a precise mathematical requirement. No dimension, semisimplicity, unitarity, or nontriviality condition is silently imposed.
\subsection{Short English statement}
Construct additional topological quantum field theories for manifolds with spin structures. The excerpt omits the known examples to be surpassed and the dimension and coefficient conventions, so it is a construction request with those choices unresolved.
\subsection{Sources}
\begin{itemize}
\item Local UnsolvedMath registry, permanent ID 10400150, source code catalogue entry, ``Problem 8.4 --- Find other spin TQFT's.'' No individual proposer is supplied.
\item Ohtsuki, \emph{Problems on invariants of knots and 3-manifolds} (2002), Problem 8.4, PDF page 123. Supplied URL: \url{https://msp.org/gtm/2002/04/gtm-2002-04-024s.pdf}; recorded access 2026-07-29 and extraction pdf-font-structure. The collection distinguishes 231 numbered bold problem, conjecture, and question headings from surrounding references.
\item The registry labels the item open and OPEN-TRIAGE and summarizes it as seeking spin TQFTs beyond known examples. The supplied status evidence is NEEDS\_REVIEW pending a current primary-literature audit. Redistribution terms for the publicly accessible author/publisher copy also require review.
\end{itemize}
\end{document}
